\chapter{声音、愿景与新星——AI 的良知、喉舌与未来力量}
\label{ch:policy-rising}
%% ============================================================
%% Chapter 14 — AI Policy, Ethics, Media Figures & Rising Stars
%% ============================================================

当技术以指数级速度演进时，
围绕技术的\textbf{叙事}——谁来讲述、谁来质疑、谁来规范——
往往决定了技术最终服务于谁。
2023--2026 年的 AI 浪潮中，
一群政策制定者、伦理学者、科普传播者、开源英雄和年轻创业者
从不同角度塑造了这场革命的面貌。
他们不一定训练模型，但他们定义了模型被训练的\emph{边界}、
被使用的\emph{方式}和被理解的\emph{语境}。

本章覆盖五个互相交织的群体：
制定规则的\textbf{政策人物}、
发出警告的\textbf{伦理学者}、
连接公众与前沿的\textbf{媒体与公共知识分子}、
构建共享基础设施的\textbf{开源社区英雄}、
以及正在重新定义可能性边界的\textbf{35 岁以下新星}。

%% ============================================================
\section{AI 政策：监管者追赶技术}
\label{sec:policy}
%% ============================================================

2023 年是 AI 监管的「觉醒之年」。
在 ChatGPT 发布不到一年后，
全球主要经济体几乎同时启动了各自的 AI 治理框架。
这种速度在技术监管史上前所未有——
相比之下，互联网经历了近 20 年才迎来第一部全面的数据隐私法规（GDPR，2018）。

\subsection{美国：行政令与芯片管制的双重策略}

美国的 AI 治理采取了一种独特的「双轨制」：
行政令（Executive Order）处理国内安全与创新，
出口管制（Export Controls）处理地缘政治竞争。

\textbf{Kamala Harris}（1964--）作为副总统，
成为 Biden 政府 AI 政策的核心推动者和公众面孔。
2023 年 10 月 30 日，Biden 签署了美国历史上第一份
AI 行政令（Executive Order on the Safe, Secure, and Trustworthy
Development and Use of Artificial Intelligence），
Harris 在签字仪式上鼓掌的画面成为标志性时刻。
该行政令要求：
开发双用途基础模型（dual-use foundation models）的公司
在发布前必须向联邦政府报告安全测试结果；
商务部制定 AI 红队测试（red-teaming）标准；
联邦机构任命首席 AI 官（Chief AI Officers）。
Harris 随后亲赴 2023 年 11 月的 Bletchley Park AI 安全峰会，
代表美国发布了关于 AI 安全的新倡议，
包括成立美国 AI 安全研究所（US AI Safety Institute，隶属 NIST）。
2024 年 3 月，她宣布了 OMB 的首份联邦 AI 使用政策，
标志着从「自愿承诺」向「制度约束」的转变。

\begin{whybox}[为什么行政令而非立法？]
美国国会在 AI 立法上进展缓慢——
截至 2026 年初，仍未通过任何联邦级别的综合性 AI 法案。
行政令虽然法律效力有限（可被下一任总统撤销），
但它能在数周内生效，不需要经历旷日持久的两党博弈。
Biden 团队选择行政令，本质上是一种
``先建立框架、再推动立法'' 的策略。
事实上，Trump 在 2025 年 1 月就任后迅速撤销了该行政令，
证明了这种路径的脆弱性。
\end{whybox}

\textbf{Bruce Reed}（1960--）是这些政策背后的关键操盘手。
作为白宫副幕僚长（Deputy Chief of Staff），
Reed 负责协调跨部门的 AI 政策制定，
被《Politico》称为 ``Biden AI 政策的大脑''。
他的团队在不到 90 天内完成了行政令的起草——
这在白宫的政策制定流程中堪称闪电速度。
Reed 此前在 Clinton 政府时期就负责过互联网政策，
他将那个时代的经验——``不要过早扼杀创新，但要建立安全护栏''——
带入了 AI 治理。

\textbf{Gina Raimondo}（1971--）作为商务部长，
掌控着美国 AI 地缘战略中最锋利的武器：芯片出口管制。
2022 年 10 月，商务部产业与安全局（BIS）发布了对华芯片出口限制，
禁止向中国出口先进 AI 芯片（如 NVIDIA A100/H100）
和先进制程半导体设备。
Raimondo 在一次 CBS 采访中坦言：
``今天的国家安全不只是坦克和导弹，而是技术。
是半导体，是 AI，是无人机。''
她主导了三轮逐步升级的出口管制：
2022 年 10 月（初始限制）、
2023 年 10 月（堵住 NVIDIA 通过降规格芯片绕过管制的漏洞）、
2024 年 12 月（扩展至 24 类芯片制造设备和 140 个实体清单新增条目）。
2025 年 1 月 13 日，在 Biden 政府的最后几天，
她发布了最全面的 AI 出口管制框架，
按国家分三级（盟友/中立/受限）分配 AI 芯片配额。
Raimondo 的策略核心是：
\emph{不追求彻底脱钩（decoupling），
而是精准管控关键瓶颈技术}。

\begin{corebox}[芯片管制的意外后果]
Raimondo 的出口管制在遏制中国获取最先进 AI 芯片方面取得了部分成功，
但也产生了深远的意外后果：
它加速了中国 AI 芯片自主化的进程
（华为昇腾 910B/C 的快速迭代即是直接回应），
催生了复杂的灰色市场走私网络，
并促使 DeepSeek 等中国实验室在软件层面进行极致优化——
DeepSeek V3 在受限的 H800 集群上实现了与 H100 集群可比的训练效率，
某种程度上证明了出口管制的``水床效应''：
按下一处，另一处鼓起。
\end{corebox}

\subsection{英国：Bletchley Park 与 AI 安全研究所}

英国选择了一条与美国和欧盟不同的道路：
不急于立法，而是聚焦于\textbf{技术安全评估}。
这一策略的执行者是
\textbf{Ian Hogarth}（1983--），
一位天使投资人和科技企业家。

Hogarth 在 2023 年 6 月发表了一篇广受关注的
\emph{Financial Times} 文章 ``We Must Slow Down the Race to God-Like AI''，
随后被 Rishi Sunak 首相任命为英国前沿 AI 任务组
（Frontier AI Taskforce）的负责人。
2023 年 11 月 1--2 日，在具有深刻象征意义的 Bletchley Park
（二战密码破解中心、Alan Turing 的工作地）
举行了全球首届 AI 安全峰会（AI Safety Summit）。
28 国代表（包括中国）签署了《Bletchley 宣言》，
承认了 frontier AI 的潜在风险，
并同意在安全测试方面开展国际合作。

峰会的核心成果是英国 AI 安全研究所
（UK AI Safety Institute, AISI）的正式成立，
Hogarth 出任首任所长。
AISI 的任务是：在 frontier model 发布前后对其进行安全评估，
检测有害能力（如生化武器知识生成、网络攻击辅助）。
这种``政府级红队测试''的模式启发了美国（NIST AI Safety Institute）、
日本和加拿大建立类似机构。

\begin{keybox}[Bletchley Park 的符号意义]
选择 Bletchley Park 作为峰会地点并非巧合：
1940 年代，Turing 和团队在这里破解了 Enigma 密码，
改变了二战进程。80 年后，
各国领导人在同一地点讨论如何应对 Turing 的智力遗产
所带来的最深远挑战。
这种历史回响赋予了峰会超越外交辞令的叙事力量。
\end{keybox}

\subsection{欧盟：AI Act——全球首部综合性 AI 法案}

如果说美国靠行政令、英国靠安全评估，
欧盟则选择了它最擅长的工具：\textbf{立法}。

EU AI Act 的立法历程始于 2021 年 4 月的欧盟委员会提案，
历经近三年的谈判，于 2024 年 3 月 13 日
以 523 票赞成、46 票反对的压倒性多数在欧洲议会通过。
这是全球第一部全面的、具有法律约束力的 AI 监管框架。

三位关键推动者值得记录：

\textbf{Thierry Breton}（1955--），
法国人，时任欧盟内部市场专员（Commissioner for Internal Market）。
Breton 是 AI Act 在欧盟委员会层面的首席推动者。
他在 AI Act 通过后高调宣布：
``欧洲现在是 AI 的全球标准制定者。
我们监管尽可能少——但必要时绝不手软。''
Breton 此前是 Atos（欧洲最大 IT 服务公司之一）的 CEO，
他将企业家的务实与监管者的决心结合在一起。
他坚持 AI Act 必须覆盖通用目的 AI 模型（GPAI），
这一条款在最终版本中成为最具争议也最具影响力的部分。

\textbf{Dragoș Tudorache}（1972--），
罗马尼亚人，欧洲议会议员，AI Act 在议会的两位联合报告人之一。
Tudorache 此前是罗马尼亚内政部长，
他回忆说自己对 AI 的兴趣始于 2015 年阅读 Nick Bostrom 的
\emph{Superintelligence} 一书。
他与意大利议员 \textbf{Brando Benifei} 搭档，
花了两年时间在欧洲议会中推动 AI Act 的通过。
Tudorache 在接受 \emph{MIT Technology Review} 采访时表示：
``没有人以前做过这样的事。
这不是从架子上拿一份旧模板就能完成的。''

AI Act 的核心是一个\textbf{风险分级框架}：
\begin{itemize}[noitemsep]
  \item \textbf{不可接受风险}（Unacceptable Risk）：
        社会信用评分、实时远程生物识别（执法例外除外）——直接禁止。
  \item \textbf{高风险}（High Risk）：
        关键基础设施、教育、就业中的 AI 系统——需通过合规评估。
  \item \textbf{有限风险}：
        聊天机器人、深度伪造——需透明度义务（标注 AI 生成内容）。
  \item \textbf{最低风险}：
        垃圾邮件过滤器等——不受额外监管。
\end{itemize}

对通用目的 AI 模型（GPAI）的特别条款要求：
所有 GPAI 提供者必须提供模型卡（model cards）和技术文档；
被认定为具有``系统性风险''的模型（训练算力超过 $10^{25}$ FLOPs）
还必须进行对抗性安全测试并向欧洲 AI 办公室报告严重事件。

\subsection{中国：渐进式、分领域的 AI 治理}

中国的 AI 治理路径与西方截然不同。
中国没有选择``一步到位的综合立法''，
而是采取了\textbf{渐进式、分领域的监管策略}——
先针对具体场景出台规定，再逐步构建统一框架。

三部关键法规构成了中国 AI 治理的基础框架：
\begin{itemize}[noitemsep]
  \item \textbf{《互联网信息服务算法推荐管理规定》}
        （2022 年 3 月生效）：
        规范推荐算法，要求算法透明度和用户选择权。
  \item \textbf{《互联网信息服务深度合成管理规定》}
        （2023 年 1 月生效）：
        针对深度伪造（deepfake）技术，
        要求内容标识和备案管理。
  \item \textbf{《生成式人工智能服务管理暂行办法》}
        （2023 年 8 月生效）：
        这是全球首个专门针对生成式 AI 的监管法规，
        要求备案审批、内容安全和数据合规。
\end{itemize}

2025 年 10 月，修订后的《网络安全法》
新增了专门的``人工智能合规''条款，
标志着 AI 治理正式被纳入中国网络安全的顶层法律架构。
同时，2025 年 8 月国务院发布的
《关于深入实施``人工智能+''行动的意见》
则代表了另一个方向：
以政策激励推动 AI 与实体经济的深度融合，
设定了到 2027 年 AI 应用普及率超 70\%、
2030 年超 90\% 的宏大目标。

\begin{whybox}[中国为什么没有出台综合性 AI 法？]
2025 年，中国曾将综合性 AI 立法列入计划议程，
但随后将其移出。
分析认为，这并非``放弃监管''，
而是务实地认识到：
在 AI 技术快速迭代的阶段，
过早的综合立法可能抑制创新。
相比之下，``先试点、再标准、再立法''的路径
允许监管者通过观察 DeepSeek、Kimi 等国产模型的实际部署效果
来校准监管力度。
这种方法更灵活，但也导致了
企业面临多套重叠规则的``法规碎片化''挑战。
\end{whybox}

\subsection{Yoshua Bengio：从学术殿堂到政策前线}

\textbf{Yoshua Bengio}（1964--）
是深度学习三巨头之一（2018 年图灵奖共同获得者），
但在 2023--2026 年间，他的角色发生了深刻转变：
从纯粹的学术研究者变成了全球 AI 安全政策的首席科学顾问。

2023 年 7 月，Bengio 在美国参议院隐私与技术小组委员会作证，
呼吁政府投资 AI 安全研究并建立监管框架。
2023 年 11 月的 Bletchley Park 峰会后，
他被委托领导《国际 AI 安全科学报告》
（International Scientific Report on the Safety of Advanced AI）——
一个受 IPCC（政府间气候变化专门委员会）启发的倡议，
由 30 多个国家的代表支持，
超过 100 位独立专家参与撰写。
临时报告于 2024 年 5 月的首尔 AI 峰会发布，
完整版于 2026 年 2 月的巴黎 AI 行动峰会发布。

Bengio 的转变在 AI 社区引发了争议。
一些同行（尤其是 Yann LeCun）批评他的安全立场过于悲观，
认为当前的 AI 系统距离 AGI 还很遥远。
Bengio 的回应是：
``我宁愿因为过早发出警告而被嘲笑，
也不愿因为沉默而在灾难发生后被追责。''

%% ============================================================
\section{AI 伦理：技术的良心}
\label{sec:ethics}
%% ============================================================

如果说政策制定者从\emph{外部}为 AI 设定边界，
伦理学者则从\emph{内部}审视 AI 的价值观、偏见和权力结构。
2020--2026 年间，一系列标志性事件——
Google 伦理 AI 团队的解散、``随机鹦鹉''论文的发表、
面部识别偏见的曝光——
将 AI 伦理从学术角落推到了公众聚光灯下。

\subsection{Google 伦理 AI 风暴：Timnit Gebru 与 Margaret Mitchell}

2020 年 12 月，Google 伦理 AI 团队联合负责人
\textbf{Timnit Gebru}（1982/1983--）的离职
成为 AI 伦理领域的分水岭事件。

Gebru 是一位厄立特里亚裔埃塞俄比亚出生的计算机科学家，
在加入 Google 前已是 AI 公平性研究的先驱——
她共同创立了 Black in AI，
一个推动更多黑人参与 AI 研发的倡导组织。
在 Google，她与 Margaret Mitchell 共同领导伦理 AI 团队，
建立了一支包含大量少数族裔成员的研究小组。

导火索是一篇论文。
Gebru 与 Emily Bender、Mitchell 等人合著了
``On the Dangers of Stochastic Parrots:
Can Language Models Be Too Big?''，
系统性地分析了大型语言模型的风险：
环境成本（训练的碳排放）、
训练数据中嵌入的社会偏见、
以及模型``看起来流利但不理解''的根本局限。
据 Gebru 所述，Google 管理层要求她撤回论文或移除所有 Google 员工的署名。
当她要求知晓审稿人身份和具体反馈时，
被告知``已经接受了辞职''——Gebru 否认自己曾辞职，
坚称这是被解雇。

两个月后，2021 年 2 月，
\textbf{Margaret Mitchell}（也被解雇。
Google 声称她违反了公司行为准则，
将文件转移到公司外部。
Mitchell——一位在自然语言生成、
视觉描述和模型偏见评估方面有深厚积累的研究者——
否认了不当行为。

\begin{corebox}[``Stochastic Parrots'' 论文的深远影响]
``随机鹦鹉''论文最初因 Google 内部的审查争议而广为人知，
但它的学术价值远超其引发的公关危机。
论文提出的核心论点——
大型语言模型本质上是在``拼接''训练数据中的模式，
而非真正``理解''语言——
成为此后所有关于 LLM 能力边界讨论的起点。
该论文发表于 2021 年 ACM FAccT 会议。
值得注意的是，Mitchell 的署名被化名为
``Shmargaret Shmitchell''，
这个黑色幽默式的伪名本身就成了对学术审查的无声抗议。
有趣的是，2026 年 3 月，
Mitchell 本人撰文 ``No, `AI' is not a Stochastic Parrot''，
对该术语的\emph{过度简化使用}进行了反思——
她区分了论文的原始论点（警告过度依赖规模）
与某些人将其简化为``AI 不可能有任何能力''的误读。
\end{corebox}

Gebru 离开 Google 后创立了
\textbf{DAIR 研究所}（Distributed AI Research Institute），
一个独立于大型科技公司的 AI 研究机构，
致力于以社区为中心的 AI 研究。
Mitchell 则于 2021 年 10 月加入 Hugging Face，
担任首席伦理科学家（Chief Ethics Scientist），
在那里她将伦理原则嵌入开源 AI 平台的架构中——
这被 \emph{TIME} 杂志 2023 年 ``AI 百大影响力人物''
选中为标志性案例，
证明伦理不只是批评，也可以是\emph{建设}。

\subsection{Emily Bender：语言学家的警告}

\textbf{Emily Bender}（约 1973--）
是华盛顿大学的计算语言学教授，
``随机鹦鹉''论文的第一作者。
与许多 AI 研究者不同，
Bender 的学术根基在\emph{语言学}而非计算机科学。
这种独特的视角使她能够提出一个
许多 AI 从业者不愿面对的问题：
\emph{LLM 是在使用语言，还是在模拟使用语言的统计模式？}

Bender 是 AI 批评界最一致的声音之一。
在 2024 年 11 月 Harvey Mudd College 的演讲中，
她将 RAG（检索增强生成）比作``用好数据做的纸浆模型——
材料更好了，但本质上还是纸浆''。
她的核心观点可以概括为三个层次：
\begin{enumerate}[noitemsep]
  \item LLM 的输出正确时，这在很大程度上是\emph{偶然}的；
  \item 流利性（fluency）不等于理解（comprehension）；
  \item 将 LLM 的输出当作``知识''是一种危险的范畴错误。
\end{enumerate}

\subsection{Joy Buolamwini：揭示算法歧视的``诗人战士''}

\textbf{Joy Buolamwini}（约 1990--）的故事
始于一个简单而令人不安的发现：
作为 MIT Media Lab 的研究生，
她发现面部识别软件无法检测到她的面孔——
除非她戴上一张白色面具。

这个个人经历推动她进行了系统性研究。
她的 \textbf{Gender Shades} 项目（2018）揭示了
IBM、Microsoft 和 Face++ 的商业面部识别系统中的严重偏见：
对深肤色女性的错误率高达 34.7\%，
而对浅肤色男性仅为 0.8\%。
这一 43 倍的差距——不是抽象的统计数字，
而是真实的人被错误识别、错误逮捕、错误排斥的风险。

Buolamwini 创立了\textbf{算法正义联盟}
（Algorithmic Justice League, AJL），
她将自己定位为``诗人战士''（poet of code），
因为她在技术抗议中融入了说唱和表演艺术。
AJL 的使命不仅是技术审计，
更是``建立受影响社区的声音和选择权''。

她的工作产生了直接的政策影响：
IBM 在 2020 年宣布退出通用面部识别市场，
明确引用了偏见问题；
多个美国城市（包括旧金山、波士顿）
通过了限制政府使用面部识别的法规。

\subsection{Meredith Whittaker：从 Google 罢工到 Signal 总裁}

\textbf{Meredith Whittaker}（约 1977--）
的职业轨迹本身就是一部 AI 伦理运动的缩影。

在 Google 工作 13 年期间，
她创立了 Google 的 Open Research 项目，
并共同创立了\textbf{AI Now 研究所}
（AI Now Institute，纽约大学）——
全球最早的跨学科 AI 社会影响研究机构之一（2017）。
2018 年 11 月，她参与组织了
Google 历史上最大规模的员工罢工（Google Walkout），
抗议公司对性骚扰投诉的处理方式以及与军方的 AI 合作（Project Maven）。

罢工后，Whittaker 表示被告知如果要留在 Google，
必须放弃 AI 伦理工作和在 AI Now 的角色。
她选择在 2019 年离开。
此后她担任 FTC（联邦贸易委员会）主席 Lina Khan 的顾问，
并于 2022 年 9 月被任命为\textbf{Signal}
（加密通信应用）的首任总裁。

Whittaker 在 Signal 的角色将她的 AI 批评
与一个具体的替代愿景结合在一起：
如果大型科技公司的商业模式本质上依赖监控，
那么 Signal 所代表的\emph{加密、非营利、开源}模式
就是一种结构性的反命题。

\subsection{Kate Crawford：AI 的``政治经济学''}

\textbf{Kate Crawford}（约 1976--）是 AI Now 研究所的联合创始人，
USC Annenberg 的研究教授。
她的 2021 年著作 \emph{Atlas of AI: Power, Politics,
and the Planetary Costs of Artificial Intelligence}
开创性地将 AI 分析的视角从``算法''扩展到``供应链''：
从内华达州的锂矿到亚马逊仓库的工人，
从 ImageNet 数据集中的标注偏见到军事监控的应用。

Crawford 的核心论点是：
AI 既不是``人工的''也不是真正``智能的''——
它是一个建立在\emph{人类劳动、自然资源提取和权力不对称}
之上的社会技术系统。
\emph{Atlas of AI} 被 \emph{New Yorker} 称为
``对机器学习训练数据的迷人历史''，
被 \emph{Science} 称为``及时而紧迫的贡献''。

\subsection{Safiya Noble：搜索引擎的种族政治}

\textbf{Safiya Umoja Noble}（约 1975--）
是 UCLA 信息学教授，
她的 2018 年著作 \emph{Algorithms of Oppression:
How Search Engines Reinforce Racism}
通过一个简单实验揭示了深层问题：
在 Google 中搜索 ``Black girls'' 得到的结果
与搜索 ``white girls'' 截然不同——
前者充斥着色情化内容，后者则相对中性。

Noble 的工作将``算法偏见''的讨论从面部识别
扩展到了\emph{信息检索}——
一个看似中立但实际上深度嵌入了社会权力结构的技术领域。
她指出，搜索引擎的排名算法不是``客观的''，
而是反映了广告商的商业利益和社会的既有偏见。

\subsection{Gary Marcus：AI 最响亮的怀疑者}

\textbf{Gary Marcus}（1970--）的角色很独特：
他既不是伦理学者也不是政策制定者，
而是一个在 AI 能力问题上最持久、最高调的\textbf{技术怀疑者}。

Marcus 是 NYU 心理学与神经科学荣休教授，
认知科学家，著有多本 AI 通俗读物。
2016 年，他创立了 Geometric Intelligence
（一家将认知科学与机器学习结合的 AI 初创公司），
后被 Uber 收购。
他还在美国参议院 AI 听证会上作证。

2022 年，Marcus 发表了引发广泛争议的文章
``Deep Learning is Hitting a Wall''，
预测深度学习在可靠性和推理能力上的根本瓶颈。
AI 圈对他的态度高度两极化——
支持者称他为``AI 的良心''，
批评者称他为``永远唱衰的人''（Elon Musk 曾转发讽刺他的 meme）。
然而，Marcus 2024 年初在 \emph{IEEE Spectrum} 的访谈中指出，
他关于可靠性问题的许多预测已经被验证：
幻觉（hallucination）仍然是 LLM 最难解决的问题之一。

Marcus 自己的总结颇为精辟：
``讽刺是针对我的；失败是深度学习之主的。''

\subsection{Melanie Mitchell：复杂性视角下的 AI}

\textbf{Melanie Mitchell}（约 1969--）
是圣塔菲研究所（Santa Fe Institute）教授，
她的独特贡献在于从\textbf{复杂系统科学}的视角分析 AI。
她的著作 \emph{Complexity: A Guided Tour}（2010，获 Phi Beta Kappa 奖）
和 \emph{Artificial Intelligence: A Guide for Thinking Humans}（2019）
为非技术读者提供了理解 AI 的深度框架。

Mitchell 的研究聚焦于``概念抽象和类比推理''——
她认为这些是人类智能的核心，
也是当前 AI 系统最薄弱的环节。
她既不属于``AI 将毁灭人类''的末日论阵营，
也不属于``深度学习能解决一切''的乐观派，
而是坚持一种基于证据的、精细的中间立场。

%% ============================================================
\section{AI 媒体与公共知识分子}
\label{sec:media}
%% ============================================================

AI 革命的独特之处在于它的\textbf{极度公开性}。
与过去的技术革命不同（半导体革命发生在洁净室里，
互联网革命发生在大学实验室里），
AI 革命很大程度上发生在 YouTube、Twitter/X、
Substack 和播客上。
一批``翻译者''——将前沿论文转化为公众理解的人——
成为了 AI 时代不可或缺的角色。

\subsection{Lex Fridman：AI 时代的 Charlie Rose}

\textbf{Lex Fridman}（1983--），
出生于苏联塔吉克斯坦的美裔计算机科学家和播客主持人。
他在 Drexel University 获得学士、硕士和博士学位，
此后在 MIT 人工智能实验室和 AgeLab 担任研究员。
2018 年起，他开始主持 \emph{Lex Fridman Podcast}，
至今已制作超过 460 期，
嘉宾涵盖 Elon Musk、Sam Altman、Mark Zuckerberg、
Jensen Huang、Yann LeCun、Jeff Bezos 乃至
印度总理 Narendra Modi。

Fridman 的访谈风格——超长形式（通常 2--4 小时）、
深度技术问题与哲学思考交织、
极其温和的追问方式——使他成为 AI 领域最有影响力的媒体人之一。
批评者认为他对嘉宾过于客气，缺乏批判性追问；
支持者则认为这种方式让嘉宾放松后更可能吐露真实想法。
无论如何，``上 Lex 的播客''已成为 AI 领袖
向科技社区传达信息的标准渠道。

Fridman 2019 年因 Elon Musk 赞扬他一篇关于 Tesla Autopilot
驾驶员注意力的研究而迅速走红——
尽管该研究未经同行评审并受到 AI 专家批评。
自 2022 年起他切换为 MIT 无薪职位，
全职投入播客和内容创作。
截至 2026 年，他的 YouTube 频道拥有数百万订阅者。

\subsection{Andrej Karpathy：从研究到教育的全栈 AI 人}

\textbf{Andrej Karpathy}（1986--）
可能是 AI 领域最成功的``知识翻译者''。
他是 OpenAI 的创始团队成员（2015），
后任 Tesla 自动驾驶（Autopilot / FSD）负责人（2017--2022），
再回到 OpenAI（2023），最终于 2024 年离开创立
\textbf{Eureka Labs}——一个``AI 原生''教育平台。

但 Karpathy 最深远的影响力来自他的\emph{教育内容}。
他在 Stanford 开设的 CS231n（卷积神经网络与视觉识别）
是深度学习教育史上最有影响力的课程之一。
他的 YouTube 频道以极其清晰的风格
从零讲解神经网络、GPT、Tokenizer 等核心概念——
他 2023 年的视频 ``Let's Build GPT from Scratch''
在数周内获得数百万播放量，
使数以万计的工程师第一次真正理解了 Transformer 的工作原理。

Karpathy 的独特之处在于他的\emph{全栈经历}：
他既做过前沿研究（ImageNet 上的 CNN），
也管理过大规模工程团队（Tesla 数百人的 Autopilot 团队），
还能将这些经验转化为任何人都能跟随的教程。
Eureka Labs 是他将这种能力制度化的尝试——
用 AI 辅助教学，让``世界上最好的老师''能够扩展到百万学生。

\subsection{Grant Sanderson (3Blue1Brown)：数学可视化的魔术师}

\textbf{Grant Sanderson}（约 1988--）
创建的 YouTube 频道 3Blue1Brown
用精美的数学动画解释了从线性代数到神经网络的核心概念。
他的``But What is a Neural Network?'' 系列（2017 年首发，2026 年更新）
可能是互联网上对神经网络最直观的入门解释。

Sanderson 的贡献看似``只是科普''，
但其实际影响不可低估：
对于一代在 AI 浪潮中成长的学生来说，
3Blue1Brown 的动画是他们理解高维空间、
反向传播和注意力机制的第一扇窗口。
2026 年 2 月，他受邀在 UC Santa Cruz 演讲，
主题是``高维球体的美与实用''——
而听众中的许多人正是因为他的视频才对数学产生了兴趣。

\subsection{Lilian Weng：从博客到 OpenAI 副总裁}

\textbf{Lilian Weng}是一个罕见的案例：
她的\emph{个人博客}可能比许多研究论文
产生了更大的教育影响力。

Weng 出生于中国，在北京航空航天大学获学士学位，
后在卡内基梅隆大学获计算机科学博士学位。
2018 年加入 OpenAI，
她从 2017 年起维护的博客 \textbf{Lil'Log}
（lilianweng.github.io）
以极其清晰、系统的长文
覆盖了从强化学习到扩散模型到
test-time compute 的几乎所有 LLM 核心主题。
每篇文章都像一部迷你教科书，
包含大量公式推导、代码片段和精心绘制的图表。

到 2025 年中，
Weng 已晋升为 OpenAI 研究与安全副总裁
（VP of Research and Safety），
接替此前一系列离职高管的空缺。
她的轨迹体现了一种独特的职业路径：
\emph{通过公开的知识贡献建立声望，
再将这种声望转化为组织内部的领导力。}

\subsection{Jeremy Howard 与 Rachel Thomas：AI 民主化的先驱}

\textbf{Jeremy Howard}（1973--）和
\textbf{Rachel Thomas}（约 1980--）
在 2016 年共同创立了 \textbf{fast.ai}，
这是 AI 教育史上最重要的项目之一。

Howard 的背景非同寻常：
他曾是 Kaggle 平台的总裁和首席科学家，
两次获得国际机器学习竞赛排名第一；
他创立了 Enlitic（第一家将深度学习应用于医疗的公司），
被 \emph{MIT Technology Review} 连续两年评为
全球 50 大最聪明公司。

fast.ai 的哲学与当时主流的``先学三年数学再碰代码''的
深度学习教育方式截然相反。
Howard 和 Thomas 借鉴了教育家 Paul Lockhart 的理念——
不要让学生在接触真正的创造之前
花费数年在干燥的理论基础上——
设计了``自上而下''的课程：
第一节课就训练一个图像分类器，
理论在实践中逐步引入。

Thomas 不仅是 fast.ai 的联合创始人，
还是 AI 伦理教育的开拓者——
她是 USF 数据伦理课程的首位讲师，
将 AI 公平性和社会影响整合进了技术教育。
她后来担任了 fast.ai 的首席执行官。

\subsection{其他重要的 AI 传播者}

\textbf{Chip Huyen}（约 1992--）是一位越南裔的 AI 系统工程师和作家。
她在越南的稻田村庄长大，
后在 Stanford 教授``机器学习系统设计''课程。
她的著作 \emph{Designing Machine Learning Systems}（2022）
成为 Amazon AI 类目的 \#1 畅销书，翻译成 10+ 种语言；
2025 年的 \emph{AI Engineering} 一书
自上市以来一直是 O'Reilly 平台上阅读量最高的书籍。
她此前在 NVIDIA（NeMo 核心开发者）、Snorkel AI 和 Netflix 工作，
还创立并出售了一家 AI 基础设施创业公司。

\textbf{Sebastian Raschka}（约 1989--）是威斯康星大学麦迪逊分校的
统计学助理教授，后加入 Lightning AI 担任首席 AI 教育官。
他的著作 \emph{Machine Learning with PyTorch and Scikit-Learn}
以及 \emph{Build a Large Language Model (From Scratch)}
以``代码优先''的风格解释了从经典 ML 到 LLM 的全栈知识。
他的博客和 Twitter 内容是 ML 社区中被引用最频繁的教育资源之一。

\textbf{李沐}（Mu Li）是前 Amazon Web Services 应用科学家，
他与 Aston Zhang 等人合著的
\emph{Dive into Deep Learning}（D2L，《动手学深度学习》）
是全球被 70 个国家 500+ 所大学采用的交互式深度学习教材。
他的中文 YouTube 频道``跟李沐学 AI''
为中文世界提供了高质量的 AI 教育内容。
D2L 的独特之处在于它是\emph{可执行的}——
每一章的代码都可以直接在 Jupyter Notebook 中运行。

\textbf{Yannic Kilcher} 是一位瑞士的 ML 研究者和 YouTube 创作者，
以逐页拆解最新 AI 论文而闻名。
他的视频覆盖了从 GPT-4 技术报告到 Mamba 论文的几乎所有重要工作，
风格直接、不加修饰，偶尔尖锐——
他是许多研究者了解新论文的第一站。

\subsection{AI 哲学家与未来学家}

\textbf{Ray Kurzweil}（1948--）
是 AI 未来学领域最具影响力（也最具争议性）的声音。
他的 2005 年著作 \emph{The Singularity Is Near}
预测``技术奇点''——人工智能超越人类智能的时刻——将在 2045 年到来。
自 2012 年起，他在 Google 担任工程总监，
致力于自然语言处理和语音识别。
2024 年，他出版了续作 \emph{The Singularity Is Nearer}，
声称 AI 的进展正在验证他二十年前的预测。
Kurzweil 的``加速回报定律''——
技术进步的速度本身在指数级加速——
对整整一代 AI 创业者产生了深刻影响，
尽管许多学术研究者认为他的预测过于乐观。

\textbf{Nick Bostrom}（1973--），
瑞典哲学家，牛津大学教授，
人类未来研究所（Future of Humanity Institute, FHI）创始主任。
他的 2014 年著作 \emph{Superintelligence: Paths, Dangers, Strategies}
是 AI 安全运动的奠基性文本之一——
从 Elon Musk 到 Bill Gates，
许多科技领袖都引用该书作为他们关注 AI 风险的原因。
Bostrom 将``超级智能''定义为
``在几乎所有认知领域大幅超越人类的智力''，
并分析了``控制问题''——
一旦超级智能出现，我们如何确保它的目标与人类一致？
FHI 于 2024 年关闭，
但 Bostrom 继续通过 Macrostrategy Research Initiative 从事 AGI 治理研究。
他 2024 年的新书 \emph{Deep Utopia} 探讨了
一个``所有问题都被 AI 解决''的世界中人类意义的问题。

\textbf{Max Tegmark}（1967--），
瑞典裔美国物理学家，MIT 教授。
他的 2017 年著作 \emph{Life 3.0: Being Human in the Age of
Artificial Intelligence} 将 AI 的讨论
放置在宇宙学的宏大框架中：
如果 Life 1.0 是生物演化（硬件和软件都由自然选择决定），
Life 2.0 是文化演化（软件可以学习，硬件仍由进化决定），
那么 Life 3.0 就是 AI 时代——
``硬件和软件都可以自我设计''。
Tegmark 共同创立了
Future of Life Institute（FLI），
该组织在 2023 年发起了要求暂停训练
比 GPT-4 更强大的 AI 模型 6 个月的公开信，
获得超过 33,000 个签名（包括 Elon Musk、Steve Wozniak、
Yoshua Bengio 等知名人物），
成为 AI 安全讨论中的标志性事件——
尽管实际上没有任何实验室真正暂停了训练。

\textbf{Yuval Noah Harari}（1976--），
以色列历史学家，
\emph{Sapiens} 和 \emph{Homo Deus} 的作者。
Harari 并非 AI 技术专家，
但他对 AI 对人类社会、劳动、意识和权力结构影响的分析
使他成为全球最受关注的 AI 评论者之一。
他在多次演讲中警告：
AI 不需要``有意识''就能颠覆人类社会——
它只需要比人类更擅长\emph{操纵叙事}。

\textbf{李开复}（Kai-Fu Lee，1961--）
在中美 AI 界都拥有独特的双重影响力。
他曾任 Apple 语音识别研究负责人、
Microsoft 全球副总裁、Google 中国区总裁，
后创立创新工场（Sinovation Ventures）。
他的 2018 年著作 \emph{AI Superpowers: China, Silicon Valley,
and the New World Order} 是向西方世界解释
中国 AI 生态系统的最有影响力的作品。
2023 年，他创立了 \textbf{01.AI}
并推出了 Yi 系列开源大模型，
使他从 AI 评论家转变为 AI 实践者——
少有人能同时在\emph{写关于} AI 和\emph{做} AI 两个维度上
保持顶级影响力。

%% ============================================================
\section{下一代：35 岁以下的新星}
\label{sec:rising}
%% ============================================================

AI 领域最令人振奋的特征之一是它的\textbf{年轻化}。
在大多数技术领域，
重大影响通常需要数十年的积累；
但在 AI 的当前浪潮中，
一群 25--35 岁的年轻人正在以创始人、
关键论文作者或核心工程师的身份
重新定义可能性的边界。

\subsection{创业新星}

\textbf{Alexandr Wang}（汪滔，1997--）
是 AI 时代``天才少年创业者''的典型代表。
他出生于新墨西哥州 Los Alamos
（他的父母都是在 Los Alamos 国家实验室工作的物理学家），
从小在科学氛围中成长。
他在 MIT 读了一年后辍学，
19 岁时（2016 年）与 Lucy Guo 共同创立了 \textbf{Scale AI}——
一家提供 AI 训练数据标注和模型评估服务的公司。

Scale AI 的核心洞察是：
AI 模型的质量瓶颈不在算法而在\emph{数据}。
公司为 OpenAI、Meta、美国国防部等客户提供高质量标注数据。
2021 年，Wang 以 24 岁的年龄
成为全球最年轻的白手起家亿万富翁。
截至 2025 年，Scale AI 估值达 290 亿美元，
\emph{Forbes} 估计 Wang 的个人净资产为 36 亿美元。

2025 年 6 月，Wang 做出了一个出人意料的决定：
离开 Scale AI，加入 Meta 担任首席 AI 官，
领导新成立的 Superintelligence Labs。
Meta 同时以 143 亿美元投资 Scale AI，获得 49\% 的股份。
这一交易使 Scale AI 估值翻倍，
也标志着 AI 行业``创始人级别人才''竞争的白热化。

\begin{keybox}[从 Los Alamos 到 Meta：一条隐喻性的路径]
Wang 的人生轨迹有一种历史性的对称：
Los Alamos 是曼哈顿计划的诞生地，
那里的科学家在 1940 年代创造了人类历史上最强大的武器；
80 年后，同一个小镇培养出的年轻人
正在构建人类历史上最强大的智能系统。
这种地理巧合提醒我们：
技术力量的集中——无论是核能还是 AI——
始终伴随着深刻的伦理责任。
\end{keybox}

\textbf{Aravind Srinivas}（约 1993--）
是 \textbf{Perplexity AI} 的联合创始人兼 CEO。
出生于印度 Chennai，
在 IIT Madras 获学士学位，
UC Berkeley 获 PhD（研究方向为语言和扩散生成模型），
期间在 OpenAI 担任研究科学家（2021--2022）。

2022 年 8 月，
他与 Denis Yarats、Johnny Ho、Andy Konwinski 共同创立了 Perplexity——
一个``AI 答案引擎''，试图重新定义搜索的范式。
与传统搜索引擎返回链接列表不同，
Perplexity 直接生成带引用来源的综合答案。
公司增长极为迅速：
截至 2025 年 9 月估值达 200 亿美元。
Srinivas 登上 2025 年 M3M Hurun India Rich List，
成为榜单上最年轻的亿万富翁。

不过，Perplexity 也面临严重的法律争议：
BBC、道琼斯和纽约时报等媒体机构指控其未经授权使用内容；
\emph{Wired} 和 Cloudflare 的分析发现，
Perplexity 使用了伪装 user-agent 的未披露爬虫，
引发了关于 AI 搜索产品与传统媒体之间利益冲突的广泛讨论。

\textbf{Scott Wu}（1997--）
是 \textbf{Cognition AI} 的联合创始人兼 CEO，
该公司开发了 Devin——被宣传为``世界第一个自主 AI 软件工程师''。
Wu 的背景极具竞争力编程色彩：
他在国际信息学奥林匹克（IOI）获得三枚金牌
（2014 年排名第一），
Codeforces 峰值 rating 达到 3350，
2021 年 Google Code Jam 第三名。
他在 Harvard 读了两年经济学后辍学，
先与他人共同创立了 AI 驱动的社交平台 Lunchclub，
再于 2023 年 11 月创立 Cognition AI。

Wu 声称 Devin 将在 2025 年底前编写公司 50\% 的代码。
无论这一目标是否实现，
Devin 的发布（2024 年 3 月）引发了
关于``AI 是否会取代程序员''的全球性讨论——
它的意义不仅在于技术能力，
更在于它迫使整个软件行业
重新思考``人类工程师''的价值定位。

\textbf{Michael Truell}（约 1997--）
是 \textbf{Anysphere} 的联合创始人兼 CEO，
该公司开发了 \textbf{Cursor}——
AI 时代增长最快的代码编辑器。
Truell 在 MIT 学习并从事研究（2017--2021），
2022 年创立 Anysphere。
Cursor 的理念是将 AI 深度集成到编码工作流中——
不是作为``聊天助手''附着在 IDE 旁边，
而是作为编辑器的\emph{原生智能层}。
截至 2026 年初，Cursor 的用户增长速度
使其成为开发者工具领域最瞩目的创业公司之一。

\textbf{Demi Guo}（约 1998--）和 \textbf{Chenlin Meng}
是 \textbf{Pika Labs} 的联合创始人。
两人在 Stanford 攻读博士期间，
参加了 Runway 举办的 AI 电影节——
对现有 AI 视频工具的失望促使她们决定自己动手。
2023 年 4 月，她们从 Stanford 退学创立 Pika。
到 2025 年 10 月，
Pika 已拥有 1450 万用户，
估值达 4.7 亿美元。
Guo 在 \emph{Fortune} 的采访中表示：
``我们不是在做专业视频工具，
而是在做一个为 Gen Z 设计的创意平台。''
Pika 的成功验证了一个重要趋势：
AI 视频生成的杀手级应用可能不是好莱坞特效，
而是\emph{社交媒体内容创作}。

\textbf{杨植麟}（Yang Zhilin，1992--），
广东汕头人，清华大学计算机系本科，
Carnegie Mellon University 计算机科学博士（导师：Ruslan Salakhutdinov 和 William Cohen）。
他是 Transformer-XL 和 XLNet 两篇里程碑论文的第一作者
（均被引用超过一万次），
曾在 Meta AI（与 Jason Weston）和 Google Brain（与 Quoc V.\ Le）工作。

2023 年，杨植麟创立了 \textbf{Moonshot AI}（月之暗面），
并推出了 Kimi 大模型。
Kimi 的差异化定位是\textbf{超长上下文}——
从一开始就支持 20 万字的输入窗口，
后扩展至百万级别——
这与他在 Transformer-XL 上的学术积累高度一致。
公司迅速成长为中国 AI 独角兽之一。
杨植麟入选 Forbes Asia 30 Under 30、NVIDIA Fellow、Siebel Scholar，
是中国``90 后''大模型创业者中学术背景最深厚的代表。

\textbf{Harrison Chase}（约 1996--）
是 \textbf{LangChain} 的联合创始人兼 CEO。
他在 Harvard 学习统计学和计算机科学，
毕业后先后在金融科技公司 Kensho（实体链接团队负责人）
和 Robust Intelligence（机器学习团队负责人）工作。

2023 年 1 月，Chase 创立 LangChain——
一个帮助开发者构建基于 LLM 的应用的开源框架。
LangChain 几乎是``AI 应用开发层''这一赛道的代名词：
它提供了链式调用（chains）、代理（agents）、
检索增强生成（RAG）和记忆管理的标准抽象。
LangChain 的 GitHub star 数在发布后数月内
飙升至数万，成为 AI 开发者最常用的中间件之一。

\subsection{学术新星}

\textbf{Tri Dao}（约 1993--）和 \textbf{Albert Gu}（约 1993--）
是两位正在重新定义深度学习底层架构的年轻研究者。

Tri Dao 的学术履历几乎完美：
Stanford 本科（数学）和博士（计算机科学，
导师 Christopher Ré 和 Stefano Ermon），
博士论文题为 ``Hardware-aware Algorithms for Efficient Machine Learning''。
他最知名的贡献是 \textbf{FlashAttention}——
一种 IO-aware 的精确注意力算法，
通过重新安排 GPU 内存层次结构（SRAM vs.\ HBM）中的计算顺序，
将 Transformer 的注意力计算加速 2--4 倍且不牺牲精度。
FlashAttention 被几乎所有主流组织采用
（Meta、Microsoft、NVIDIA、OpenAI、Google、DeepSeek 等），
集成到了所有主要 ML 框架
（PyTorch、JAX、Hugging Face transformers、vLLM、DeepSpeed 等）。

Dao 于 2023 年共同创立了 \textbf{Together AI}
（担任首席科学家），
2024 年加入 Princeton 大学担任助理教授——
同时在产业界和学术界保持活跃。

Albert Gu 此前在 Stanford 和 CMU 工作，
是结构化状态空间模型（Structured State Space Models, S4）
的核心提出者。
2023 年 12 月，
Gu 和 Dao 联合发表了 \textbf{Mamba} 论文——
一个基于选择性状态空间模型的序列建模架构，
实现了\emph{线性时间复杂度}的序列建模，
首次在语言模型中展示了与 Transformer 可比的性能。
Mamba 论文引发了 AI 社区的强烈反响——
这是自 2017 年 ``Attention Is All You Need'' 以来，
第一个被认真讨论为``Transformer 的潜在替代者''的架构。

\begin{whybox}[为什么 Mamba 如此重要？]
Transformer 的自注意力机制具有 $O(n^2)$ 的序列长度复杂度，
这意味着处理长序列的成本以二次方增长。
Mamba 通过选择性状态空间实现了 $O(n)$ 复杂度，
同时保持了内容感知推理（content-based reasoning）的能力——
这是此前所有亚二次架构的致命弱点。
虽然截至 2026 年，Transformer 仍然是主流架构，
但 Mamba 及其后续版本（Mamba-2 基于``状态空间对偶性''理论，
揭示了 SSM 和注意力变体之间的深层数学联系）
正在开辟一个 hybrid 架构的新方向。
\end{whybox}

\textbf{Hyung Won Chung}（约 1989--）
是推理模型研究领域的核心人物之一。
他在 MIT 获得博士学位后先加入 Google Brain（2019--2023），
在那里与 Jason Wei 长期合作，
共同作为第一作者发表了
``Scaling Instruction-Finetuned Language Models''（Flan-T5/PaLM 论文），
这是指令微调（instruction tuning）领域最具影响力的工作之一。

2023 年，他加入 OpenAI 的 ChatGPT 团队，
成为 \textbf{o1 推理模型}的核心贡献者之一。
o1 是 OpenAI 在 2024 年 9 月发布的专用推理模型，
标志着``让模型学会思考''的新范式。
2025 年，Chung 离开 OpenAI 加入 Meta（后为 Meta Superintelligence Labs）。
他离开后公开分享了自己的深刻思考：
``AI 是历史上最强大的杠杆——超越人类劳动、资本和代码。''

\textbf{谢赛宁}（Xie Saining，约 1992--）
是计算机视觉领域最具影响力的年轻研究者之一。
他是 NYU Courant 研究所的助理教授，
此前在 Meta AI（FAIR）工作了四年多，
2024--2025 年在 Google DeepMind 的 GenAI 团队。

他与合作者共同开发了多个广泛采用的架构：
\textbf{ResNeXt}（深度残差网络的扩展）、
\textbf{ConvNeXt}（证明了经过精心设计的 CNN
在 Transformer 统治视觉任务的时代仍然极具竞争力）、
以及 \textbf{DiT}（Diffusion Transformer，
将 Transformer 架构引入扩散模型图像生成，
被认为是 Sora 等视频生成模型的架构基础之一）。
他还是 MoCo v2（对比学习）和 MAE（掩码自编码器）等
里程碑工作的核心作者。

2026 年初，谢赛宁共同创立了 \textbf{AMI Labs}，
担任首席科学官（CSO），
致力于开发超越纯语言系统的``世界模型''——
能够直接从连续的、真实世界的感官输入中学习。
他入选了 \emph{MIT Technology Review} 的 Innovators Under 35。

\textbf{Lianmin Zheng}（郑连民）
是高效 LLM 推理领域的关键贡献者。
作为 UC Berkeley 和 Stanford 的研究者，
他是 \textbf{vLLM} 推理引擎
和 \textbf{SGLang}（结构化语言模型程序高效执行系统）
的核心作者。

vLLM 通过 PagedAttention 技术
重新设计了 KV cache 的内存管理，
实现了远超此前系统的吞吐量——
截至 2026 年，vLLM 在 GitHub 上拥有超过 73,000 stars，
成为 LLM 推理部署的事实标准之一。
SGLang 则进一步引入了 RadixAttention
（KV cache 复用）和压缩有限状态机
（加速结构化输出解码）等创新，
发表于 NeurIPS 2024。

\textbf{稚晖君 / 彭志辉}（Peng Zhihui，1993--）
是中国 AI 社区中最具传奇色彩的人物之一。
他毕业于电子科技大学（UESTC），
获信息与通信工程硕士学位（2018），
在校期间就在 Bilibili 上通过 DIY 技术项目
（自平衡机器人、微型电视、机械臂``Dummy''）
积累了大量粉丝，被称为``野生钢铁侠''。

他通过华为的``天才少年''计划
（年薪百万级别的顶尖毕业生项目）加入华为，
2023 年离开创立了 \textbf{Agibot}（智元机器人），
一家总部位于上海的具身智能和人形机器人创业公司。
Agibot 的机器人可以执行推轮椅、汽车制造等多种任务，
一款机器人可举起 40 公斤。
公司获得了比亚迪、高瓴资本等的投资，
截至 2025 年 3 月估值达 21 亿美元。

2025 年 4 月，彭志辉作为最年轻的代表（32 岁）
参加了李强总理主持的企业家座谈会——
这在中国的政商关系语境中是一个极具象征意义的信号，
表明具身智能已被提升到国家战略层面。

%% ============================================================
\section{开源社区英雄}
\label{sec:opensource}
%% ============================================================

AI 革命不仅是大公司的游戏。
一个由开源贡献者、独立研究者和社区维护者组成的生态系统
在很大程度上决定了 AI 技术的\emph{可及性}。
他们不一定出现在融资新闻里，
但他们的代码被数百万人使用。
而在 2025--2026 年间，
一个又一个``一个人改变整个行业''的故事反复上演——
从保加利亚的 C 语言极客到奥地利的 PDF 工匠，
开源社区英雄们证明了
在 AI 时代，\emph{个体的杠杆率}达到了史无前例的高度。

\subsection{Peter Steinberger 与 OpenClaw：从 PDF 工匠到 GitHub 之王}

2025 年 11 月 24 日的晚上，
\textbf{Peter Steinberger}（约 1985--）——
一位居住在维也纳和伦敦之间的奥地利软件工程师——
坐下来做了一个周末实验：
将 WhatsApp 消息接口连接到 Anthropic 的 Claude API，
让 AI 代理能够通过聊天应用执行任务。
这个实验花了大约一个小时。
他将它命名为 \textbf{Clawdbot}——
``Claude'' 与 ``claw''（龙虾爪）的文字游戏，
配上一只卡通龙虾作为吉祥物。

接下来发生的事情，
堪称开源史上最不可思议的增长曲线。

Steinberger 的背景与典型的 AI 创业者截然不同。
他是维也纳理工大学（Vienna University of Technology）校友，
17 岁时就在 iOS 开发社区崭露头角。
2011 年，他创立了 \textbf{PSPDFKit}——
一个处理 PDF 渲染和标注的移动端框架，
最终被安装在超过 10 亿台设备上，
客户包括 Dropbox、Autodesk、IBM 等巨头。
2021 年，Insight Partners 对 PSPDFKit 进行了
1 亿欧元的战略投资，
Steinberger 随后退出日常管理。
用他自己的话说：
``我花了 13 年做 PDF 工具，
赚了很多钱，
但也几乎燃尽了自己。''

退出 PSPDFKit 后，
Steinberger 转向 Web 开发和 AI 工具探索。
Clawdbot 最初只是他众多实验中的一个——
一个开源的、本地部署的自主 AI 代理框架，
允许用户通过 WhatsApp、Telegram 等消息应用
与 AI 助手互动，
AI 可以自主执行代码、管理文件、调用 API。
项目使用 TypeScript 和 Swift 编写，
以 MIT 许可证发布。

然后，它爆发了。

2026 年 1 月初，Clawdbot 在 Hacker News 和 Twitter/X 上引爆。
三天内获得 60,000 个 GitHub star。
十天内达到 210,000。
Steinberger 一个人在 80 天内提交了超过 6,600 次 commit，
编写了 518,000 行代码——
平均每天 6,475 行，
是一个资深工程师正常产出的 65 倍。
秘诀？他同时运行多个 AI 代理辅助编码，
创造了一种``人类指挥、AI 执行''的极致开发模式。

名字经历了戏剧性的变迁：
2026 年 1 月 26 日，
因 Anthropic 提出商标顾虑（``Clawdbot'' 太像 ``Claude''），
项目更名为 \textbf{Moltbot}；
三天后的 1 月 29 日，
在完成商标调查后最终定名为 \textbf{OpenClaw}。
更名期间还催生了 \textbf{Moltbook}——
一个由 140 万个 AI 代理互相对话、
人类只能观看的``AI 社交网络''，
成为 2026 年初最超现实的互联网景观之一。

2026 年 2 月 12 日，
Steinberger 登上了 \textbf{Lex Fridman Podcast \#491}，
在近四小时的访谈中讲述了 OpenClaw 的创造故事。
他在节目中坦言：
``有人出价数十亿美元，我不在乎。
我已经有过一次大退出，发现它让我空虚。''
\textbf{Sam Altman}（OpenAI CEO）
和 \textbf{Mark Zuckerberg}（Meta CEO）
都亲自联系了他。

2026 年 2 月 14 日——情人节——
Steinberger 在个人博客宣布加入 \textbf{OpenAI}，
负责``将 AI 代理带给每一个人''。
同时宣布 OpenClaw 将移交给一个独立的开源基金会，
``保持开放和独立''。
值得注意的是，
他每月自掏腰包约 20,000 美元支付 API 费用
来维护这个开源项目——
这个事实本身就成为了关于
开源 AI 可持续性的重要讨论素材。

2026 年 3 月 3 日，
OpenClaw 以超过 250,000 stars
超越 React（243,000 stars），
成为 GitHub 上 star 数最多的软件项目——
React 花了十年建立的记录，
OpenClaw 在四个月内打破。
截至 2026 年 3 月中旬，star 数已突破 325,000。

\begin{corebox}[一个人、一小时、一个时代]
Steinberger 的故事之所以如此震撼，
不仅因为增长速度的绝对数字，
更因为它揭示了 AI 时代一个深刻的结构性变化：
\emph{个体开发者的杠杆率已经达到了历史最高点}。
一个人，利用 AI 辅助编码，
可以在数周内构建出
此前需要数十人团队花费数年才能完成的软件。
OpenClaw 不仅是一个开源项目——
它是``vibe coding''（直觉编程）
和 AI 代理时代的\emph{宣言}。
它也引发了严肃的安全讨论：
一个拥有自主执行能力的 AI 代理框架，
以 MIT 许可证向所有人开放，
没有内置的安全护栏——
这是民主化，还是潘多拉的盒子？
\end{corebox}

\subsection{Clément Delangue 与 Hugging Face：开源 AI 的``GitHub''}

\textbf{Clément Delangue}（约 1986--），
法国人，Sciences Po（巴黎政治学院）毕业，
主修政治科学和经济学——
一个看似与 AI 无关的背景。
他 17 岁时就成为法国 eBay 排名最高的卖家之一，
展现了早期的创业天赋。

2016 年，Delangue 与 Julien Chaumond 和 Thomas Wolf
共同创立了 Hugging Face——
最初是一个有趣的 AI 聊天伴侣项目。
在发现 NLP 社区急需一个开放的模型共享平台后，
他们转型为模型仓库和开发者平台。

截至 2026 年，Hugging Face 已成为 AI 领域的
\emph{事实标准基础设施}：
超过 100 万个模型在平台上托管，
超过 5,000 家组织使用其工具。
Hugging Face 的 \texttt{transformers} 库
是世界上被引用和使用最多的 ML 代码库之一。
平台估值超过 45 亿美元。

Delangue 的核心哲学是``开放''：
不是作为商业策略，而是作为\emph{信念}——
AI 的发展应该是协作的、可访问的、社区驱动的。
他经常将 Hugging Face 类比为``AI 的 GitHub''，
这个类比在 2026 年 2 月变得更加贴切——
llama.cpp 和 GGML 的创始团队正式加入了 Hugging Face。

\subsection{Georgi Gerganov：让 AI 在笔记本电脑上运行}

\textbf{Georgi Gerganov}，保加利亚索菲亚的软件工程师，
是``本地 AI 运动''的奠基者。
他的背景出人意料——
Gerganov 拥有医学物理硕士学位，
此前在 ViewRay（一家放射治疗技术公司）担任首席科学家，
从事 MRI 引导放疗系统的软件开发。
这种对底层性能优化的痴迷
最终将他引向了一个完全不同的领域。

2022 年 9 月，在 OpenAI 发布 Whisper 语音识别模型后不久，
Gerganov 创建了 \textbf{whisper.cpp}——
Whisper 的纯 C/C++ 移植版本，
无需 Python 和 PyTorch 即可在消费级硬件上运行语音转文字。
这个项目验证了一个关键直觉：
\emph{许多 AI 模型并不真正需要 GPU 集群，
只要推理代码足够精炼}。

在此基础上，他创建了 \textbf{GGML}——
一个极简的 C/C++ 机器学习张量库
（核心代码仅数个源文件），
专注于 Transformer 推理的极致效率。

2023 年 3 月，在 Meta 发布 LLaMA 权重后不到一周，
Gerganov 发布了 \textbf{llama.cpp}——
一个纯 C/C++ 实现的 LLaMA 推理引擎，
能够在\emph{没有 GPU 的普通笔记本电脑}上运行大型语言模型。
这一项目从根本上改变了``本地 AI''的可能性边界。

llama.cpp 的技术创新包括：
\begin{itemize}[noitemsep]
  \item 高效的量化推理
        （支持 2-bit 到 8-bit 的多种量化格式）；
  \item 纯 CPU 推理的极致优化
        （利用 SIMD、内存映射等底层技术）；
  \item \textbf{GGUF} 格式——
        成为本地 AI 模型分发的事实标准。
\end{itemize}

截至 2026 年 2 月，llama.cpp 在 GitHub 上拥有超过 95,000 stars。
Gerganov 的前期投资者包括 Nat Friedman（前 GitHub CEO）
和 Daniel Gross。
2026 年 2 月 20 日，Gerganov 和他的团队
（包括 Xuan-Son Nguyen 和 Aleksander Grygier）
正式加入 Hugging Face——
这意味着开源 AI 生态系统的``推理层''（llama.cpp）
和``分发层''（Hugging Face Hub）
以及``模型定义层''（transformers）
第一次合并在一个组织下。
没有其他组织拥有这样的全栈控制力。

\begin{corebox}[一个人改变一个行业]
Gerganov 的案例挑战了``AI 需要数十亿美元投资''的叙事。
他以一个人的力量（后来扩展到小团队），
用不到 5 个源文件的核心代码库，
使数百万人能够在自己的设备上运行 AI 模型。
在一个由数十亿参数模型和数万块 GPU 集群主导的行业中，
llama.cpp 证明了精巧的工程
可以让``小''变得和``大''一样有影响力。
\end{corebox}

\subsection{Justine Tunney：Llamafile 与``一键运行 AI''}

\textbf{Justine Tunney} 是一位独立开发者、系统程序员和安全研究者，
在底层系统编程领域拥有近乎传奇的声望。
她此前在 Google 担任高级软件工程师（2013--2018），
以创建 \textbf{Cosmopolitan Libc}——
一个允许 C 程序编译一次、
在 Linux/macOS/Windows/FreeBSD/OpenBSD/NetBSD
六个操作系统上无修改运行的跨平台 C 运行时——
和 \textbf{Redbean}（单文件 Web 服务器）闻名。
她也是 Occupy Wall Street 运动中
\#OccupyWallSt.org 的早期技术志愿者，
这种``技术服务于社会运动''的信念贯穿了她的职业生涯。

2023--2024 年，在 Mozilla 的资助下，
Tunney 将 Cosmopolitan Libc 的跨平台魔法
应用到了 AI 推理领域，
创建了 \textbf{llamafile}——
一个将大型语言模型打包为\emph{单个可执行文件}的项目。
llamafile 基于 llama.cpp，
但做了一件 llama.cpp 没有做的事：
将模型权重和推理引擎\emph{融合}为一个文件。
用户可以双击一个文件就运行 AI 模型，
不需要安装 Python、配置 CUDA 或理解命令行。
更重要的是，Tunney 在 CPU 推理层面
实现了显著的性能突破——
她发现 llama.cpp 的矩阵乘法内核
没有充分利用现代 CPU 的 SIMD 指令集，
通过手写汇编和底层优化，
在某些硬件上将推理速度提升了数倍。

llamafile 是``AI 民主化''的极端案例：
它将使用 AI 的门槛降低到了``能双击文件''。
在一个 AI 基础设施日益复杂的时代，
这种``返璞归真''的工程哲学具有独特的价值——
它提醒我们，\emph{简单就是可及性的终极形式}。

\subsection{TheBloke (Tom Jobbins)：量化之神}

在 2023--2024 年的本地 AI 社区中，
\textbf{TheBloke}（真名 Tom Jobbins）
是一个几乎神话级别的存在。

他在 Hugging Face 上拥有超过 26,000 名关注者，
上传了\emph{数千个}量化模型——
几乎每当一个新的开源模型发布，
TheBloke 就会在数小时内提供
GGUF、GPTQ 和 AWQ 等多种量化格式的版本。
对于本地 AI 社区的用户来说，
``去 TheBloke 的页面找模型''
曾经是每天的例行操作。

2024 年 1 月后，TheBloke 停止了活跃更新，
社区迅速感受到了这个空缺——
Reddit 的 r/LocalLLaMA 上出现了大量
``谁能像 TheBloke 一样做量化？''的帖子。
一个人的贡献之大，以至于他的缺席
成为了社区讨论的主题。

\subsection{Stella Biderman 与 EleutherAI：开源 AI 运动的火种}

\textbf{Stella Biderman} 是 \textbf{EleutherAI} 的执行总监——
这个从 Discord 服务器起步的草根研究集体
可以说\emph{启动}了现代开源 AI 运动。

Biderman 是一位数学家和 AI 研究者
（芝加哥大学数学学士，Georgia Tech 计算机科学硕士），
此前在 Booz Allen Hamilton 工作。
EleutherAI 成立于 2020 年，
最初是对 OpenAI 宣布 GPT-3 但不公开权重的\emph{回应}——
一群志愿者决定自己复现一个开放的 GPT 级别模型。

他们的成果改变了开源 AI 的格局：
\textbf{GPT-Neo}（2021，GPT-3 的首个开源替代品）、
\textbf{GPT-NeoX-20B}（2022，当时最大的开源 LLM 之一）、
\textbf{Pythia}（系统化的 LLM scaling 实验套件）、
以及参与 \textbf{BLOOM}（BigScience 176B 参数模型）。

Biderman 的使命声明很明确：
``确保学术界、审计机构和非营利研究者
能够研究大规模 AI 技术，
而不需要与出售这些技术的公司合作。''
EleutherAI 从一个草根集体发展为正式的非营利研究机构，
Biderman 在其中的角色从技术贡献者转变为
整个开源 AI 生态系统的关键组织者。

\subsection{Teknium 与 Nous Research：微调的艺术}

\textbf{Teknium}（网络化名）是 \textbf{Nous Research} 的核心人物。
Nous Research 定位为``美国开源 AI 运动的领导者''，
专注于开源语言模型的训练和微调。

Teknium 在 GitHub 和 Twitter 上以其名活跃，
他的早期项目 GPTeacher（使用 GPT-4 生成模块化数据集
用于指令微调、角色扮演和代码训练）
展示了一种新的研究范式：
用合成数据（synthetic data）大规模提升开源模型的能力。

Nous Research 的 Hermes 系列模型
以及他们在合成数据生成、推理训练和
agent 架构方面的工作，
使其成为开源社区中最受尊重的微调团队之一。
他们的工作哲学是：
``不可限制的可用性和使用，
以及对科学和公众理解的促进。''

\subsection{Simon Willison：AI 工具的``瑞士军刀''匠人}

\textbf{Simon Willison}（约 1979--），
英国人，Django Web 框架的联合创始人。
在 AI 时代，他通过两个项目获得了新一轮的广泛影响力。

\textbf{Datasette}——一个将 SQLite 数据库
转化为即时可查询 API 和 Web 界面的工具——
看似与 AI 无关，但它体现了 Willison 的核心理念：
数据应该是可访问的、可探索的、可组合的。

\textbf{LLM}——一个 CLI 工具和 Python 库，
让用户从命令行与 OpenAI、Anthropic、Google、Meta 等
数十种大型语言模型交互，
将提示和响应存储在 SQLite 中，
支持嵌入生成、结构化内容提取和工具执行。
2025 年 5 月发布的 LLM 0.26
引入了工具支持（tool use），
允许 LLM 直接在终端中执行 Python 函数。

但 Willison 最独特的贡献可能是他的\emph{博客}
（simonwillison.net）。
他以极其勤勉的写作习惯，
几乎每天发布关于 AI 工具、安全隐患和实践技巧的深度文章。
他的博客是许多开发者追踪 AI 工具生态发展的首选信息源——
不是因为他做了最前沿的研究，
而是因为他做了最好的\emph{翻译、评估和整合}。

\subsection{Ollama：让本地 AI 变成``一行命令''}

如果说 llama.cpp 证明了 AI 可以在笔记本上运行，
\textbf{Ollama} 则让这件事变得\emph{任何人都能做到}。

Ollama 由 \textbf{Jeffrey Morgan} 和 \textbf{Michael Chiang}
于 2023 年创立，总部位于 Palo Alto。
两人曾参加 Y Combinator 2021 冬季批次（当时创立的是 Infra，
一个身份管理工具）。
Morgan 的早期职业生涯与 Docker 生态深度绑定——
他是 \textbf{Kitematic}（后被 Docker 收购）的创始人，
曾在 Docker 担任高级工程师和产品经理，
更早时在 Twitter 和 Google 实习。
这种``把复杂基础设施包装成简单体验''的产品直觉，
直接塑造了 Ollama 的设计哲学。

Ollama 的核心创新不在算法层面——
它底层依赖 llama.cpp——
而在\emph{用户体验}层面：
一条命令（\texttt{ollama run llama3}）
就能自动下载模型、配置推理参数、启动本地 API 服务。
它借鉴了 Docker 的镜像管理模式：
\texttt{ollama pull}、\texttt{ollama list}、\texttt{ollama run}——
对任何用过容器的开发者来说都是零学习曲线。

截至 2026 年 3 月，Ollama 在 GitHub 上拥有超过 166,000 stars，
支持 Llama、Qwen、DeepSeek、Gemma、GLM 等数十个模型家族。
它已成为本地 AI 生态的\emph{事实标准入口}——
Open WebUI、OpenClaw 等众多项目
都以 Ollama 作为默认的本地模型后端。

\begin{keybox}[Docker 思维的 AI 化]
Ollama 的成功揭示了一个重要模式：
AI 基础设施的瓶颈往往不在\emph{算法}，
而在\emph{打包和分发}。
就像 Docker 通过容器化让``在我机器上能跑''
变成了``在任何机器上都能跑''，
Ollama 通过极简的 CLI 抽象层
让``本地运行 AI 模型''
从极客的专利变成了所有开发者的日常工具。
Morgan 的 Docker 背景在这里产生了直接的知识迁移。
\end{keybox}

\subsection{Timothy Jaeryang Baek 与 Open WebUI：本地 AI 的前端革命}

有了 Ollama 做后端，
社区仍然需要一个``好看、好用、功能丰富''的前端界面。
\textbf{Open WebUI}（最初名为 Ollama WebUI）
填补了这个空白。

创始人 \textbf{Timothy Jaeryang Baek}（Tim Baek）
毕业于伦敦大学（一等荣誉学位），
后在 Simon Fraser University 获计算机科学硕士学位。
他同时也是 Mozilla Builders 社区成员。
Open WebUI 最初是 Baek 为自己的本地 AI 需求
开发的 Web 界面，
但很快因其功能的完整性和设计的精致度
在社区中获得了爆炸性增长。

Open WebUI 提供的功能列表
几乎涵盖了一个本地 AI 助手可能需要的一切：
RAG（检索增强生成）、离线支持、多模型切换、
图像生成集成、语音交互、
插件系统和 Markdown 渲染。
它完全自托管，可以在不联网的情况下运行——
对隐私敏感的用户和企业来说，这是核心卖点。

截至 2026 年 3 月，
Open WebUI 在 GitHub 上拥有超过 126,000 stars，
成为 Ollama 生态中最受欢迎的前端。
``Ollama + Open WebUI''的组合
被社区公认为搭建本地 AI 助手的``黄金搭档''——
类似于早年 LAMP（Linux + Apache + MySQL + PHP）
作为 Web 开发标准栈的地位。

\subsection{AUTOMATIC1111：Stable Diffusion 的``人民界面''}

在文本生成领域之外，
\textbf{图像生成}的开源民主化同样经历了
由社区英雄驱动的关键时刻。

2022 年 8 月，Stability AI 开源了 Stable Diffusion，
但原始代码需要命令行操作和 Python 编程知识。
一位匿名开发者 \textbf{AUTOMATIC1111}
（真实身份从未公开）在 Stable Diffusion 发布后数周内
创建了 \textbf{Stable Diffusion WebUI}——
一个基于 Gradio 的 Web 界面，
将图像生成的全部参数暴露为直观的滑块和输入框。

这个项目的影响力是变革性的。
它支持 txt2img（文本生成图像）、
img2img（图像到图像转换）、
Inpainting（局部重绘）、ControlNet 集成、
LoRA 加载、Textual Inversion、
面部修复、超分辨率放大等数十项功能。
更关键的是，它建立了一个\emph{扩展生态系统}——
社区开发者可以编写插件扩展 WebUI 的能力，
形成了一个自组织的创新网络。

截至 2026 年 3 月，
AUTOMATIC1111 的 stable-diffusion-webui
在 GitHub 上拥有超过 162,000 stars，
是除 OpenClaw 外 star 数最多的 AI 项目之一。
匿名身份使 AUTOMATIC1111 成为了一个
纯粹以\emph{作品}定义影响力的案例——
没有创始人故事，没有 VC 融资，没有公关团队，
只有代码和社区。

\subsection{ComfyUI 与 Comfyanonymous：节点化 AI 创作的先锋}

如果说 AUTOMATIC1111 的 WebUI 是
``让每个人都能生成图像''，
\textbf{ComfyUI} 则是
``让高级用户完全控制生成流程''。

2022 年 10 月，一位匿名开发者 \textbf{comfyanonymous}
在发现 Stable Diffusion 后被深深吸引——
据项目官方叙述，最初的动机
是为了``生成更好的兽耳角色图像''。
这个看似轻浮的起点
催生了 AI 图像生成领域最强大的开源工具之一。

ComfyUI 的核心创新是\textbf{节点化工作流}（node-based workflow）：
用户通过连接可视化节点来构建图像生成管线，
每个节点代表一个操作
（如加载模型、添加噪声、应用 ControlNet、执行采样）。
这种设计灵感来自专业视效软件
（如 Nuke、Houdini、Blender 的 Shader Editor），
给予用户对扩散过程每一个步骤的精确控制。

ComfyUI 从个人项目发展为正式公司 \textbf{Comfy Org}，
由 \textbf{Yoland Yan} 担任 CEO。
2025 年 9 月，Comfy Org 完成了 1,700 万美元融资，
目标是成为``创意 AI 的操作系统''。
截至 2026 年 3 月，
ComfyUI 在 GitHub 上拥有超过 106,000 stars，
已支持图像、视频、3D 资产和音频生成，
被无数艺术家、工作室和开发者用于创意生产。

\begin{whybox}[为什么两个 Stable Diffusion UI 可以共存？]
AUTOMATIC1111 和 ComfyUI 看似竞争，
实际上服务不同的用户群体和使用场景。
AUTOMATIC1111 的 WebUI 类似 Photoshop——
功能完整、即开即用、学习曲线平缓；
ComfyUI 类似 After Effects 或 Houdini——
节点化、可组合、学习曲线陡峭但上限极高。
前者服务``想要生成图像的人''，
后者服务``想要精确控制生成过程的人''。
这种分层共存模式在开源生态中反复出现——
就像 VS Code 和 Vim 的关系。
\end{whybox}

\subsection{Kohya：LoRA 训练的无名英雄}

在 Stable Diffusion 生态中，
``生成图像''只是故事的一半；
另一半是``训练自己的模型''。
日本开发者 \textbf{kohya-ss}（匿名）
编写的 \textbf{sd-scripts} 训练脚本
成为了社区进行 \textbf{LoRA}
（Low-Rank Adaptation）微调的事实标准工具。

LoRA 是一种参数高效的微调方法，
允许用户用少量图像（通常 20--50 张）
和有限的 GPU 资源
训练出能够生成特定角色、风格或概念的模型适配器。
kohya-ss 的脚本使这一流程\emph{可操作化}——
从数据准备到训练参数配置到模型导出，
覆盖了完整的工作流。
社区开发者 \textbf{bmaltais} 为这些脚本
构建了 Gradio GUI 界面（kohya\_ss GUI），
进一步降低了使用门槛。

kohya-ss 的脚本在 GitHub 上拥有约 7,000 stars
（bmaltais 的 GUI 封装约 12,000 stars），
数字看似不起眼，
但它们的实际影响力远超 star 数所能反映的范围——
Hugging Face 和 Civitai 上数以万计的
LoRA 微调模型都是用这些脚本训练的。
kohya-ss 是``基础设施级''的开源贡献者：
不显眼，但移除它，整个生态会塌陷。

\subsection{vLLM 与 Inferact：从学术论文到推理基础设施}

\textbf{vLLM} 的故事始于一个简单的观察：
LLM 推理的瓶颈不在计算，
而在\emph{内存管理}。

2023 年，UC Berkeley 的博士生
\textbf{Woosuk Kwon}、\textbf{Zhuohan Li}、
\textbf{Siyuan Zhuang}、\textbf{Lianmin Zheng}
以及 Stanford 的 \textbf{Ying Sheng}
等人发表了 \textbf{PagedAttention} 论文
（发表于 SOSP 2023），
提出了一个受操作系统虚拟内存和分页技术启发的
注意力算法：
将 KV cache 存储在非连续的内存块中，
通过``页表''（page table）进行管理，
几乎消除了内存碎片和冗余复制。
在此基础上构建的 vLLM 推理引擎
实现了接近零浪费的 KV cache 管理
和灵活的跨请求 KV cache 共享，
吞吐量比当时的主流方案高 2--4 倍。

vLLM 迅速成为\emph{生产级 LLM 部署}的首选引擎，
被 Anyscale、Databricks、多家 LLM API 提供商采用。
2025 年，vLLM 的核心维护者——
\textbf{Simon Mo}、\textbf{Woosuk Kwon}、
\textbf{Kaichao You}、\textbf{Roger Wang}——
将项目商业化为创业公司 \textbf{Inferact}。
2026 年 1 月，Inferact 宣布完成
1.5 亿美元种子轮融资（估值 8 亿美元），
由 Andreessen Horowitz 和 Lightspeed 联合领投。
这是 AI 推理基础设施领域最大的种子轮之一，
标志着行业关注点正从\emph{训练}向\emph{推理}转移。

\begin{keybox}[从训练时代到推理时代]
vLLM/Inferact 的融资是一个
行业范式转移的信号。
2023--2024 年，资本主要流向``训练更大的模型''
（基础模型公司、GPU 集群）；
2025--2026 年，瓶颈转向了``高效地服务这些模型''——
延迟、吞吐量、冷启动时间、成本效率
成为了新的竞争维度。
PagedAttention 从一篇学术论文
到一家 8 亿美元估值公司的路径，
也证明了基础设施层的``无聊创新''
往往比应用层的``炫酷产品''
拥有更持久的商业价值。
\end{keybox}

%% ============================================================
\section{交叉与张力}
\label{sec:tensions}
%% ============================================================

本章覆盖的人物看似分布在截然不同的领域——
政策、伦理、媒体、开源、创业——
但他们之间存在深刻的\textbf{结构性张力}：

\begin{itemize}[noitemsep]
  \item \textbf{开放 vs.\ 安全}：
        开源社区（EleutherAI、Hugging Face、Nous Research、OpenClaw）
        坚信透明和可访问性是安全的\emph{前提}；
        而部分安全倡导者（Bostrom 的某些论点、
        英国 AISI 的评估框架）则担忧
        开放权重模型和自主代理框架可能被恶意使用。
        OpenClaw 的 MIT 许可证发布方式
        将这一张力推到了新的高度——
        一个``无护栏''的自主 AI 代理框架向所有人开放。
  \item \textbf{监管 vs.\ 创新}：
        EU AI Act 的分级框架在保护公民权利的同时，
        被批评者指责增加了 AI 创业者的合规成本，
        可能使欧洲在 AI 竞赛中进一步落后。
  \item \textbf{批评 vs.\ 建设}：
        Gebru 和 Bender 的批评被一些人视为必要的警告，
        被另一些人视为对技术进步的阻碍。
        Mitchell 从``被解雇的批评者''到``Hugging Face 首席伦理科学家''
        的转变，暗示了一种可能的综合：
        伦理不必只是``说不''，也可以是``更好地说是''。
  \item \textbf{年轻 vs.\ 经验}：
        25 岁的创始人（Wang、Wu、Chase）
        正在构建价值数十亿美元的公司，
        而 70 多岁的哲学家（Kurzweil、Bostrom）
        仍在塑造关于 AI 未来的根本叙事。
        这种代际共存——
        实践者和思考者、建设者和批评者——
        正是当前 AI 生态系统活力的来源。
\end{itemize}

\begin{whybox}[为什么``谁''的问题和``什么''的问题一样重要？]
本书前几章讨论了 AI 的技术架构、训练方法和部署策略。
但技术从来不是在真空中发展的。
\emph{谁}获得资金、\emph{谁}制定规则、\emph{谁}发出警告、
\emph{谁}教育下一代——
这些``谁''的问题在很大程度上决定了
AI 技术最终是增进还是削弱人类福祉。

Timnit Gebru 被 Google 解雇的故事
不仅仅是一个雇佣纠纷——
它是关于\emph{谁有权定义 AI 的风险}的权力斗争。
Alexandr Wang 从 Los Alamos 到 Meta 的轨迹
不仅仅是一个创业成功故事——
它是关于\emph{AI 产业权力如何集中}的案例研究。
Georgi Gerganov 的 llama.cpp
不仅仅是一个技术项目——
它是关于\emph{AI 可及性}的政治宣言。

理解这些人的故事，
就是理解 AI 革命的深层结构。
\end{whybox}
